\begin{figure}[t]
  \centering
  \begin{tikzpicture}[
    node distance=1.8cm and 2.2cm,
    every node/.style={
      rectangle,
      draw=black,
      line width=1.2pt,
      text width=1.8cm,
      text centered,
      font=\small,
      minimum height=1.4cm
    },
    arrow/.style={
      ->,
      line width=1.2pt,
      >= stealth
    }
  ]

    % Define colors (RGB to normalize for TikZ)
    \colorlet{garnet}{Garnet}
    \colorlet{atlantic}{Atlantic}
    \colorlet{congaree}{Congaree}
    \colorlet{lightbg}{black!10}

    % ===== INPUT STAGE =====
    % Radar sensor
    \node[
      fill=garnet,
      text=white,
      line width=1.5pt,
      draw=black
    ] (radar) at (0, 2) {X-band\\Marine Radar};

    % Camera sensor
    \node[
      fill=garnet,
      text=white,
      line width=1.5pt,
      draw=black
    ] (camera) at (0, 0) {Camera\\(RGB)};

    % ===== DATA GENERATION =====
    % PPI image output
    \node[
      fill=lightbg,
      draw=black,
      line width=1.2pt
    ] (ppi) at (2.2, 2) {PPI\\Image};

    % RGB image output
    \node[
      fill=lightbg,
      draw=black,
      line width=1.2pt
    ] (rgb) at (2.2, 0) {RGB\\Image};

    % ===== CALIBRATION STAGE =====
    \node[
      fill=atlantic,
      text=white,
      line width=1.5pt,
      draw=black
    ] (calibration) at (4.4, 1) {Calibration\\(\textit{intrinsic},\\extrinsic)};

    % ===== PROJECTION STAGE =====
    \node[
      fill=congaree,
      text=white,
      line width=1.5pt,
      draw=black,
      text width=2.0cm
    ] (projection) at (6.6, 1) {Radar-to-Camera\\Projection};

    % ===== LABEL GENERATION =====
    \node[
      fill=atlantic,
      text=white,
      line width=1.5pt,
      draw=black
    ] (bboxgen) at (8.8, 1) {Bounding Box\\Generation};

    % ===== NOISY LABEL TRAINING =====
    \node[
      fill=congaree,
      text=white,
      line width=1.5pt,
      draw=black,
      text width=2.0cm
    ] (noisytraining) at (11.0, 1) {Noisy Label\\Training\\Pipeline};

    % ===== OUTPUT STAGE =====
    \node[
      fill=garnet,
      text=white,
      line width=1.5pt,
      draw=black,
      text width=1.8cm
    ] (detector) at (13.2, 1) {Trained\\Camera\\Detector};

    % ===== ARROWS =====
    % Radar to PPI
    \draw[arrow, black] (radar) -- (ppi);

    % Camera to RGB
    \draw[arrow, black] (camera) -- (rgb);

    % PPI and RGB to Calibration
    \draw[arrow, black] (ppi) -- (calibration);
    \draw[arrow, black] (rgb) -- (calibration);

    % Calibration to Projection
    \draw[arrow, black] (calibration) -- (projection);

    % Projection to BBox
    \draw[arrow, black] (projection) -- (bboxgen);

    % BBox to Noisy Training
    \draw[arrow, black] (bboxgen) -- (noisytraining);

    % Noisy Training to Detector
    \draw[arrow, black] (noisytraining) -- (detector);

    % ===== FEEDBACK LOOP (optional, visual indicator) =====
    % Arc showing feedback from detector output to noisy training
    \draw[
      arrow,
      black,
      dashed,
      line width=1pt
    ] (detector.south) -- ++(0, -0.8) -- ++(0, 0)
      node[below=0.3cm, text=black, font=\tiny] {}
      -- ++(-2.2, 0) -- ++(0, 0.8) -- (noisytraining.south);

  \end{tikzpicture}
  \caption{High-level pipeline overview for radar-guided camera labeling. The pipeline integrates X-band marine radar and RGB camera data through calibration and projection steps to generate noisy training labels for a camera-based object detector.}
  \label{fig:pipeline_overview}
\end{figure}
